\chapter[Apresentação da 2ª edição]{Apresentação da 2ª edição}

Passados dez meses de implantação da atual Política de Segurança em
Tecnologia da Informação e Comunicação (PSTIC), novos recursos e fluxos
tecnológicos foram implantados, tornando necessária sua atualização.

O desafio de instituir uma política institucional dentro de uma
autarquia pública federal que, durante muitos anos, atuou em processos
diversos e difusos sem padronização e formalidade, coloca-nos em
situação de maior atenção e responsabilidade.

Identificamos resistências e dúvidas quanto aos procedimentos
institucionais e maior necessidade de elucidação das ações.

O desafio de garantir a governança preconizada para a administração
pública envolve três funções básicas, alinhadas às tarefas sugeridas
pela ISO/IEC 38500:2008:

\begin{enumerate}
\item avaliar o ambiente, os cenários, o desempenho e os resultados atuais
    e futuros;

\item direcionar e orientar a preparação, a articulação e a coordenação de
    políticas e planos, alinhando as funções organizacionais às necessidades
    das partes interessadas (usuária/os dos serviços, cidadãos e sociedade
    em geral) e assegurando o alcance dos objetivos estabelecidos;

\item monitorar os resultados, o desempenho e o cumprimento de políticas e
    planos, confrontando-os com as metas estabelecidas e as expectativas das
    partes interessadas.
\end{enumerate}

O Acórdão 1768/2022 do TCU visa contribuir para a transformação digital
do país, conscientizando os gestores públicos acerca dos riscos aos
quais as organizações estão sujeitas, de modo que sejam implementados
controles e medidas de segurança adequados para enfrentá-los.

Assim, o TCU estabelece os ``Cinco controles de segurança cibernética
para ontem'' para embasar a adoção, pelas organizações públicas
federais, dos controles críticos de cibersegurança que se seguem:

\begin{itemize}[]
\item \textbf{Controle 1 --- Inventário e controle de ativos corporativos}:
    identificar e impedir a utilização de ativos de TI não
    autorizados/gerenciados como vetores de ataques cibernéticos;

\item \textbf{Controle 2 --- Inventário e controle de ativos de
      \emph{software}}: identificar e impedir a utilização de \emph{softwares}
    não autorizados/gerenciados como vetores de ataques cibernéticos;

\item \textbf{Controle 7 --- Gestão contínua de vulnerabilidades}:
    evitar a exploração de vulnerabilidades conhecidas nos ativos
    corporativos de TI;

\item \textbf{Controle 14 --- Conscientização sobre segurança e
      treinamento de competências}: reduzir a possibilidade de incidentes e
    ataques derivados do comportamento humano (engenharia social);

\item \textbf{Controle 17 --- Gestão de respostas a incidentes}: melhorar a
    capacidade de identificar potenciais ameaças e ataques, evitar que se
    espalhem e recuperar rapidamente dados e sistemas eventualmente
    corrompidos.
\end{itemize}

Também a publicação da Portaria CRP-06 nº~99/2024, que ``estabelece a
jornada de trabalho de 30 horas semanais e o teletrabalho no âmbito do
Conselho Regional de Psicologia da 6ª Região, cria e regulamenta o
regime híbrido de trabalho e dá outras providências'' nos coloca a
necessidade de regulamentar com mais rigor o uso das tecnologias da
informação e comunicação.

Esta nova edição apresenta inovações em diferentes âmbitos, seja pela
inserção de tutoriais de acesso aos sistemas e programas, seja pela
atualização de ferramentas para abertura de chamados e reformulação de
fluxos;

propõe uma ponte com o Código de Ética da Sociedade Brasileira de
Computação, um avanço na vinculação prática do compromisso ético de
nossa profissão aos princípios ali indicados para a oferta de serviços
de qualidade a sociedades;

incorpora normativas convergentes na defesa intransigente da democracia
e da Psicologia: a Lei de Acesso à Informação (LAI; Lei nº~12.527/2011),
o Marco Civil da Internet (Lei nº~12.965/2014) e a Lei Geral de Proteção
de Dados (LGPD; Lei nº~13.709/2018), fundamentando, assim a expressão
prática do espírito destas.

Assim, apresentamos a segunda edição da presente Política com o foco no
aprimoramento, na qualificação e na inovação dos fluxos e procedimentos
institucionais.

Talita Fabiano de Carvalho

Conselheira presidenta do CRP-06